\documentclass[10pt,letterpaper]{article}
\usepackage[T1]{fontenc}
\usepackage[utf8]{inputenc}
\usepackage{lmodern}
\usepackage[margin=0.76in]{geometry}
\usepackage[table]{xcolor}
\usepackage{array,tabularx,longtable,booktabs}
\usepackage{enumitem}
\usepackage{amssymb,amsmath}
\usepackage{hyperref}
\usepackage{fancyhdr,lastpage}
\usepackage{microtype}
\usepackage{parskip}
\definecolor{cdblue}{HTML}{203957}
\definecolor{cdgray}{HTML}{586677}
\hypersetup{colorlinks=true,linkcolor=cdblue,urlcolor=cdblue,pdftitle={Delivery Engineering Technical Reference}}
\pagestyle{fancy}
\fancyhf{}
\lhead{\small\color{cdgray}Continuous Delivery Engineering}
\rhead{\small\color{cdgray}Technical reference}
\cfoot{\thepage\ / \pageref{LastPage}}
\setlength{\headheight}{14pt}
\setlength{\emergencystretch}{2em}
\setlist[itemize]{leftmargin=*,itemsep=3pt,topsep=4pt}
\setlist[enumerate]{leftmargin=*,itemsep=3pt,topsep=4pt}
\begin{document}
\begin{center}
{\LARGE\bfseries\color{cdblue}Delivery Engineering Technical Reference}\\[7pt]
{\large Release capability}\\[6pt]
{\small Reference edition: 2026-10-30 $\bullet$ Original educational and operational material}
\end{center}
\noindent\color{cdblue}\rule{\linewidth}{1.2pt}\color{black}
\vspace{0.4em}
\section{Purpose and normative status}
This document is an original engineering reference. Statements using "must" or "shall" define recommended controls for adoption after local risk review, not requirements imposed by the Humble--Farley text.


\section{Terminology and semantics}
\begin{itemize}
\item Continuous integration builds and tests regularly integrated changes; delivery aims to keep changes deployable; deployment executes release to an environment; feature release exposes capability to users.
\item Define artifact, candidate, environment, gate, approval, provenance, change set, blast radius, SLO and release train in the project glossary.
\end{itemize}
\subsection{Verification}
\begin{itemize}
\item Name the accountable owner and expected evidence for each assertion.
\item Demonstrate a negative case or rollback condition, not only the happy path.
\item Record an explicit exception if a control cannot be met.
\end{itemize}


\section{Quality attributes}
\begin{itemize}
\item Reproducibility: identical tracked inputs and controlled toolchain produce equivalent verifiable outputs.
\item Recoverability: documented recovery actions have rehearsed evidence and safe data handling.
\item Auditability: every deployment maps to source SHA, artifact digest, approver, gates and runtime version.
\end{itemize}
\subsection{Verification}
\begin{itemize}
\item Name the accountable owner and expected evidence for each assertion.
\item Demonstrate a negative case or rollback condition, not only the happy path.
\item Record an explicit exception if a control cannot be met.
\end{itemize}


\section{Value stream}
\begin{itemize}
\item Map idea, code, review, CI, package, security, acceptance, staging, decision, production, observation.
\item For every stage record owner, wait time, execution time, failure probability, required evidence and queue policy.
\end{itemize}
\subsection{Verification}
\begin{itemize}
\item Name the accountable owner and expected evidence for each assertion.
\item Demonstrate a negative case or rollback condition, not only the happy path.
\item Record an explicit exception if a control cannot be met.
\end{itemize}


\section{Responsibility}
\begin{itemize}
\item Application owner owns code and runbooks; platform owner owns shared guardrails and identity boundaries; security owner defines enforceable policy; operations owns monitoring and incident coordination.
\item A gate must have a named escalation path; avoid hidden operational dependencies.
\end{itemize}
\subsection{Verification}
\begin{itemize}
\item Name the accountable owner and expected evidence for each assertion.
\item Demonstrate a negative case or rollback condition, not only the happy path.
\item Record an explicit exception if a control cannot be met.
\end{itemize}


\section{Minimum evidence contract}
\begin{itemize}
\item Decision/requirement identifier and current version.
\item Executable configuration or code reference at an immutable revision.
\item Artifact digest or migration ID, pipeline run URL, environment and timestamp.
\item Gate outcome, reviewer identity, exception record if applicable.
\item Post-change observation and recovery evidence.
\end{itemize}

\section{Review prompts}
\begin{itemize}
\item What would invalidate this model under a concurrent deployment?
\item Which checks apply before publication versus at deployment time?
\item What data or side effects can survive rollback?
\item How would an auditor reproduce the decision six months later?
\end{itemize}

\section{Source attribution}
Book foundation: Humble and Farley, Continuous Delivery (2010), especially Chapters 1--15 as indicated by topic. Modern implementation detail is supplementary and should be checked against current vendor documentation listed in the suite README.



\end{document}
